\section*{Bab \ref{ch:b2:equilibrium-shifts} --- Tetapan Kesetimbangan dan Pergeseran Kesetimbangan}

\begin{solution}{exo:b2:equilibrium-shifts:1}
$K^\circ = \exp(32\,800/(8.314 \times 298.15)) = \num{5.6e5}$, dan untuk
$+\qty{32.8}{kJ/mol}$ kebalikannya, \num{1.8e-6}. Faktor 100 memerlukan
$\Delta(\Delta_r G^\circ) = -RT\ln 100 = \qty{-11.4}{kJ/mol}$.
\end{solution}

\begin{solution}{exo:b2:equilibrium-shifts:2}
$\ln K^\circ(400) = \ln(5.4 \times 10^5) + (91\,880/8.314)(1/400 -
1/298.15) = 13.20 - 9.44 = 3.76$, $K^\circ \approx 43$. Tabel memberikan
35.6: dalam rentang \qty{100}{K}, pendekatan $\Delta_r H^\circ$ tetap sudah
meleset 20\,\% (reaksi menjadi lebih \omterm{def:b2:reaction-enthalpy:exothermic}{eksoterm} ketika $T$ naik).
\end{solution}

\begin{solution}{exo:b2:equilibrium-shifts:3}
(a) $n = 1$, $r = 0$, $\varphi = 2$: $v = 1$ (tekanan uap merupakan fungsi
$T$). (b) $v = 3 - 1 + 2 - 3 = 1$. (c) $v = 3 - 1 + 2 - 1 = 3$. (d) Empat
spesies, satu reaksi, dua fase (karbon dan gas): $v = 4 - 1 + 2 - 2 =
3$; jika gasnya hanya dibuat dari karbon dan uap air, $n(\ce{CO}) = n(\ce{H2})$
dalam gas, $s = 1$ dan $v = 2$.
\end{solution}

\begin{solution}{exo:b2:equilibrium-shifts:4}
Pemanasan: ke arah \omterm{def:b2:reaction-enthalpy:exothermic}{endoterm}, menuju \ce{NO2} (ke arah maju). Pemampatan:
menuju molekul gas yang lebih sedikit, \ce{N2O4} (ke arah balik). Argon
pada volume tetap: tidak ada perubahan, tekanan parsial gas-gas yang
bereaksi tidak berubah.
\end{solution}

\begin{solution}{exo:b2:equilibrium-shifts:5}
Pada \qty{298}{K}: $\Delta_r G^\circ = 180.6 - 298.15 \times 0.0248 =
\qty{173.2}{kJ/mol}$, $K^\circ = \num{4.5e-31}$. Pada \qty{2000}{K}: $\Delta_r
G^\circ = 180.6 - 2000 \times 0.0248 = \qty{131.0}{kJ/mol}$, $K^\circ =
\num{3.8e-4}$. Dengan $\Delta\nu_{\text{gas}} = 0$, $x(\ce{NO})^2 = K^\circ x(\ce{N2})
x(\ce{O2})$: $x(\ce{NO}) = \sqrt{3.8 \times 10^{-4} \times 0.78 \times 0.21} =
\num{7.9e-3}$, hampir satu persen. Mesin membakar bahan bakar pada suhu
setinggi itu; ketika gas mendingin, $K^\circ$ turun sangat jauh, tetapi
penguraian \ce{NO} sangat lambat: kesetimbangan suhu tinggi itu
``membeku'' di dalam gas buang (sampai sebuah katalis
menghilangkannya).
\end{solution}

\begin{solution}{exo:b2:equilibrium-shifts:6}
Dengan $\xi$ mol ester dan jumlah zat total yang tetap, $Q =
\xi^2/[(1 - \xi)(n_B - \xi)]$. Untuk $n_B = 1.0$: $\xi^2/(1 - \xi)^2 = 4$,
$\xi
= \qty{0.67}{mol}$. Untuk $n_B = 3.0$: $3\xi^2 - 16\xi + 12 = 0$,
$\xi =
\qty{0.90}{mol}$. Pengambilan air menjaga $Q$ tetap di bawah
$K^\circ$ berapa pun $\xi$: reaksi terus berjalan ke arah maju sampai
asamnya habis.
\end{solution}

\begin{solution}{exo:b2:equilibrium-shifts:7}
Tanpa argon: jumlah zat $1 - \alpha$, $2\alpha$, total $1 + \alpha$;
$Q = 4\alpha^2/(1 - \alpha^2) = 0.148$, $\alpha = \sqrt{0.148/4.148} = 0.189$.
Dengan \qty{1}{mol} argon pada \qty{1}{bar}: total $2 + \alpha$, $Q =
4\alpha^2/[(1 - \alpha)(2 + \alpha)] = 0.148$, sehingga $4.148\alpha^2 + 0.148\alpha
- 0.296 = 0$, $\alpha = 0.250$: pengenceran pada tekanan total tetap
menurunkan tekanan parsial, yang menguntungkan sisi dengan molekul gas
lebih banyak. Pada volume tetap, tekanan parsial gas-gas yang bereaksi
tidak berubah: $\alpha$ tetap 0.189.
\end{solution}

\begin{solution}{exo:b2:equilibrium-shifts:8}
Hanya padatan: $n = 3$, $r = 1$, $s = 1$ ($p(\ce{NH3}) = p(\ce{HCl})$),
$\varphi = 2$: $v = 3 - 1 - 1 + 2 - 2 = 1$; pilihlah $T$, maka tekanannya
sudah tertentu (sebuah ``tekanan disosiasi''). Dengan penambahan amonia,
$s = 0$ dan $v = 2$: $T$ dan satu tekanan dapat dipilih; hasil kali
$p(\ce{NH3})\,p(\ce{HCl})$ ditetapkan oleh $T$.
\end{solution}

\begin{solution}{exo:b2:equilibrium-shifts:9}
$\Delta_r H^\circ = -R\ln(K_2/K_1)/(1/T_2 - 1/T_1) = -8.314 \times
\ln(0.0173)/(1/600 - 1/500) = \qty{-101}{kJ/mol}$, lebih negatif daripada
pada \qty{298}{K} (\qty{-91.9}{kJ/mol}), seperti yang diramalkan oleh
\omterm{thm:b2:reaction-enthalpy:kirchhoff}{hukum Kirchhoff} ($\Delta_r C_p^\circ < 0$).
\end{solution}

\begin{solution}{exo:b2:equilibrium-shifts:10}
$Q = \dfrac{n_{\ce{NH3}}^2\,n_{\text{tot}}^2}{n_{\ce{N2}}n_{\ce{H2}}^3}
\left(\dfrac{p^\circ}{p}\right)^2$. Pada $T$, $p$ tetap:
$\partial\ln Q/\partial n_{\ce{N2}} = -1/n_{\ce{N2}} + 2/n_{\text{tot}}$,
yang positif jika $n_{\ce{N2}} > n_{\text{tot}}/2$, yaitu
$x(\ce{N2}) > 1/2$. Maka $Q$ melampaui $K^\circ$ dan kesetimbangan
bergeser ke arah balik. Pada volume tetap,
$\partial\ln Q/\partial n_{\ce{N2}} = -1/n_{\ce{N2}} < 0$: selalu ke arah
maju.
\end{solution}

\begin{solution}{exo:b2:equilibrium-shifts:11}
$K^\circ = \exp(-1620/(8.314 \times 1100)) = 0.84$: $p(\ce{CO2}) =
\qty{0.84}{bar}$ pada kesetimbangan, yaitu $n = pV/RT = 0.84 \times 10^5 \times
0.0100/(8.314 \times 1100) = \qty{0.092}{mol}$ gas. Dengan \qty{0.050}{mol}
batu kapur, seluruhnya terurai ($p = \qty{0.46}{bar} < K^\circ p^\circ$):
tidak ada kesetimbangan, kesetimbangan terputus. Dengan \qty{0.200}{mol},
terjadi kesetimbangan: \qty{0.092}{mol} \ce{CO2} dan \ce{CaO}, tersisa
\qty{0.108}{mol} \ce{CaCO3}.
\end{solution}

\begin{solution}{exo:b2:equilibrium-shifts:12}
Unggun 1 keluar pada sekitar \qty{890}{K} dan konversi 68\,\%, unggun 2
pada sekitar \qty{785}{K} dan 92\,\%, unggun 3 pada sekitar \qty{725}{K}
dan 98\,\%. Dalam satu unggun adiabatik, gas memanas selama bereaksi
sampai mencapai kurva kesetimbangan, yang pada suhu itu hanya
memungkinkan sekitar 68\,\%. Katalis yang dingin sepanjang waktu bukanlah
pilihan karena katalis hanya bekerja dalam keadaan panas; pendinginan di
antara unggun mempertahankan laju dan menaikkan batas kesetimbangan.
Unggun terakhir dibatasi oleh kesetimbangan pada suhu terendah yang masih
memungkinkan katalis bekerja (pabrik kemudian menyerap trioksidanya dan
melewatkan gasnya sekali lagi melalui katalis).
\end{solution}

\begin{solution}{pb:b2:equilibrium-shifts:1}
\textbf{1.} $\Delta_r H^\circ = \qty{-91.88}{kJ/mol}$, $\Delta_r S^\circ =
2(192.77) - 191.61 - 3(130.68) = \qty{-198.1}{J/(K.mol)}$, $\Delta_r G^\circ =
\qty{-32.8}{kJ/mol}$, $K^\circ = \num{5.6e5}$.
\textbf{2.} $K^\circ(700) = \exp(-54\,380/(8.314 \times 700)) = \num{8.8e-5}$;
$K^\circ(800) = \exp(-77\,320/(8.314 \times 800)) = \num{8.9e-6}$.
\textbf{3.} $\Delta_r H^\circ = -R\ln(K_{800}/K_{700})/(1/800 - 1/700)$,
yaitu \qty{-106}{kJ/mol}, lebih negatif daripada pada \qty{298}{K}.
\textbf{4.} $\ln K^\circ(723) = \ln(8.75 \times 10^{-5}) + 12\,777 \times
(1/723.15 - 1/700)$, dengan $12\,777 = 106\,230/8.314$: $K^\circ =
\num{4.9e-5}$.
\textbf{5.} $\Delta_r G^\circ(723) = -91.88 + 723.15 \times 0.1981 =
\qty{51.4}{kJ/mol}$, $K^\circ = \num{1.9e-4}$: empat kali terlalu besar,
karena $\Delta_r H^\circ$ dan $\Delta_r S^\circ$ berubah dalam rentang
\qty{425}{K} ($\Delta_r C_p^\circ \approx \qty{-44}{J/(K.mol)}$).
\textbf{6.} \ce{N2} $1 - \xi$, \ce{H2} $3 - 3\xi$, \ce{NH3} $2\xi$, total $4 -
2\xi$.
\textbf{7.} $x = 2\xi/(4 - 2\xi)$, sehingga $1 - x = (4 - 4\xi)/(4 - 2\xi)$;
fraksi nitrogen $(1 - \xi)/(4 - 2\xi) = (1 - x)/4$, dan fraksi hidrogen
tiga kali lebih besar.
\textbf{8.}
\[
  K^\circ = \frac{x^2}{\frac{1-x}{4}\left(\frac{3(1-x)}{4}\right)^3}
  \left(\frac{p^\circ}{p}\right)^2 = \frac{256}{27}\,\frac{x^2}{(1 - x)^4}
  \left(\frac{p^\circ}{p}\right)^2 ;
\]
lalu ambillah akar kuadratnya.
\textbf{9.} $a = (\sqrt{27}/16) \times \sqrt{4.9 \times 10^{-5}} \times 1 =
\num{2.27e-3}$; $x/(1 - x)^2 = a$ memberikan $x = 0.0023$.
\textbf{10.} $a = 0.455$; $ax^2 - (2a + 1)x + a = 0$: $x = 0.25$.
\textbf{11.} $n = 3$, $r = 1$, $s = 1$ (perbandingan 1 : 3 dipertahankan
oleh reaksi), $\varphi = 1$: $v = 2$. Begitu $T$ dan $p$ dipilih, tidak ada
lagi yang bebas.
\textbf{12.} Tekanan: ke arah maju ($\Delta\nu_{\text{gas}} = -2$). Suhu:
ke arah balik (\omterm{def:b2:reaction-enthalpy:exothermic}{eksoterm}).
\textbf{13.} Gas inert pada tekanan total tetap menurunkan tekanan parsial
gas-gas yang bereaksi, seperti halnya penurunan tekanan: ke arah balik.
\textbf{14.} $\partial\ln Q/\partial n_{\ce{N2}} = (-1/0.30 + 2)/n_{\text{tot}} < 0$:
$Q$ turun di bawah $K^\circ$, ke arah maju.
\textbf{15.} Tidak: pada volume tetap $\partial\ln Q/\partial n_{\ce{N2}} =
-1/n_{\ce{N2}} < 0$ apa pun komposisinya: ke arah maju.
\textbf{16.} Karena hampir tidak ada amonia, $Q$ jauh di bawah $K^\circ$ di
saluran masuk: gas masuk jauh dari kesetimbangan dan bereaksi ke arah
maju.
\textbf{17.} Gas terlalu singkat berada di atas katalis sehingga
kesetimbangan tidak tercapai; kontak yang lebih lama menuntut konverter
yang lebih besar dan lebih mahal demi keuntungan yang kecil.
\textbf{18.} $x = 2\xi/(4 - 2\xi) = 0.18$ memberikan $\xi = \qty{0.305}{mol}$
per mol \ce{N2} yang diumpankan: \omterm{def:b2:equilibrium-shifts:single-pass}{konversi sekali lewat} hidrogen (dan
nitrogen) $0.305$, yaitu 31\,\%.
\textbf{19.} Dalam operasi tunak, seluruh umpan segar pada akhirnya
terkonversi: \qty{2.00}{mol} amonia keluar per satuan waktu; konverter
harus dilewati $1.00/
0.305 = \qty{3.3}{mol}$ \ce{N2} per satuan waktu,
sisanya didaur ulang.
\textbf{20.} Argon (dan metana) yang masuk bersama umpan tidak pernah
bereaksi dan tidak ikut terbuang bersama amonia cair: tanpa \omterm{def:b2:equilibrium-shifts:single-pass}{purge}, keduanya
akan menumpuk di dalam sirkuit dan menurunkan tekanan parsial pereaksi.
\textbf{21.} Pada \qty{500}{K} katalis besi terlalu lambat: kesetimbangannya
memang menguntungkan, tetapi tidak akan pernah didekati dalam waktu yang
layak secara industri.
\textbf{22.} Pada \qty{723}{K} dan \qty{200}{bar}, dengan umpan
stoikiometris dan gas sempurna: $\boldsymbol{x(\ce{NH3}) \approx 0.25}$.
\end{solution}
